\section{Population selection as an optional bridge}

This appendix records one explicit population-genetic construction for the boundary introduced in
Chapter~\ref{sec:book-protein-scales}.  It is not needed for the operator or MSA arguments.  Its role
is to show exactly what must be added before an assay endpoint can receive a population-dynamic
interpretation.

On the finite sequence space $\mathcal X=[q]^L$, let $n_x(t)$ be the expected abundance of genotype
$x$.  Let $Q_{\mathrm{mut}}$ be an irreducible mutation generator in the column convention:
$Q_{xy}\geq0$ is the rate from $y$ to $x$ for $x\neq y$, and $\sum_xQ_{xy}=0$.  For an infinite
asexual population in a fixed environment, with frequency-independent Malthusian rates $m_c(x)$ and
no recombination,
\begin{equation}
\frac{dn}{dt}=\mathcal L_cn,
\qquad
\mathcal L_c=Q_{\mathrm{mut}}+\operatorname{diag}(m_c).
\label{eq:book-mutation-selection-operator}
\end{equation}
After normalization by total abundance, $p=n/(\mathbf1^{\mathsf T}n)$ obeys the
replicator--mutator equation
\begin{equation}
\frac{dp}{dt}=Q_{\mathrm{mut}}p+
\left[\operatorname{diag}(m_c)-\overline m_c(p)I\right]p,
\qquad
\overline m_c(p)=\sum_xm_c(x)p_x.
\label{eq:book-replicator-mutator-equation}
\end{equation}
The irreducible Metzler operator $\mathcal L_c$ has a positive principal right eigenvector,
\begin{equation}
\mathcal L_cv_c=\lambda_cv_c,
\qquad
p_c^\star=\frac{v_c}{\mathbf1^{\mathsf T}v_c},
\qquad
\overline m_c(p_c^\star)=\lambda_c.
\label{eq:book-mutation-selection-perron-equilibrium}
\end{equation}
This equilibrium is an eigenproblem coupling mutation and growth.  It follows from the declared
population model, not from exponentiating an arbitrary assay value.

A different, reversible stationary construction gives a Gibbs-like law.  Suppose a coarse-grained
substitution chain has rates $K_c(x\mid y)$ from $y$ to $x$, the neutral chain has full-support
stationary law $\pi_{\mathrm{mut}}$, and a dimensionless selection potential $\Psi_c$ is derived
from a fixation or diffusion model such that
\begin{equation}
\frac{K_c(x\mid y)}{K_c(y\mid x)}
=\frac{\pi_{\mathrm{mut}}(x)}{\pi_{\mathrm{mut}}(y)}
\exp\!\left[\Psi_c(x)-\Psi_c(y)\right].
\label{eq:book-selection-local-detailed-balance}
\end{equation}
Detailed balance then gives
\begin{equation}
\pi_c(x)=
\frac{\pi_{\mathrm{mut}}(x)e^{\Psi_c(x)}}
{\sum_{y\in\mathcal X}\pi_{\mathrm{mut}}(y)e^{\Psi_c(y)}},
\qquad
H_c(x)=-\log\pi_{\mathrm{mut}}(x)-\Psi_c(x)+\log Z_c.
\label{eq:book-reversible-selection-law}
\end{equation}
The potential $\Psi_c$ is a coordinate of this specified construction.  It can be identified with
$Y_e$, $m_c$, or a model score only after a separate derivation and dimensional calibration.

Finite-population drift, recombination, changing environments, frequency dependence, and genealogy
alter the state dynamics or the observation law.  They belong to the population model and should not
be silently absorbed into one scalar called fitness.
